\documentclass[11pt,a4paper]{article}
\usepackage[utf8]{inputenc}
\usepackage{amsmath,amssymb,amsthm}
\usepackage{listings}
\usepackage{geometry}
\usepackage{hyperref}
\geometry{margin=2.5cm}

\lstset{
  language=C++,
  basicstyle=\ttfamily\small,
  keywordstyle=\bfseries,
  breaklines=true,
  frame=single,
  numbers=left,
  numberstyle=\tiny,
  tabsize=4,
  showstringspaces=false
}

\title{IOI 2005 -- Riv (Rivers)}
\author{Solution}
\date{}

\begin{document}
\maketitle

\section{Problem Statement Summary}
There are $N$ villages arranged in a tree (rooted at village 0, which has a lumber camp). Each village $v$ (for $v = 1, \ldots, N-1$) has:
\begin{itemize}
    \item $w_v$: amount of wood produced,
    \item $d_v$: distance to its parent in the tree.
\end{itemize}

We need to build exactly $K$ sawmills (at villages, not at the root which already has one). The cost of transporting wood from village $v$ is $w_v \times (\text{distance from } v \text{ to the nearest sawmill on path to root})$. Minimize the total transportation cost.

\section{Solution Approach}

\subsection{DP Formulation}
We define:
\[
dp[v][j][u] = \text{min cost for subtree of } v \text{, using } j \text{ sawmills, where } u \text{ is the nearest ancestor with a sawmill}
\]

However, tracking the exact ancestor $u$ is expensive. Instead, we note that what matters is only the \emph{distance} from $v$ to the nearest ancestor sawmill. We can track this by passing the nearest ancestor sawmill's identity along the DFS.

\textbf{Cleaner formulation:} During DFS, we know the set of ancestors. The nearest ancestor with a sawmill determines the transport cost. Since we process top-down, we can pass the distance to the nearest ancestor sawmill as a parameter.

Define:
\[
dp[v][j] = \text{min cost for subtree of } v \text{, using } j \text{ sawmills in the subtree (including possibly at } v\text{)}
\]
parameterized by the distance $D$ from $v$ to its nearest ancestor sawmill.

Since $D$ can take at most $O(\text{depth})$ distinct values (one for each ancestor of $v$), and the depth is at most $N$, we can enumerate over ancestors.

\subsection{Implementation}
For each vertex $v$, iterate over all ancestors $u$ of $v$ (the potential ``nearest ancestor sawmill''). For each such $u$, compute the cost. Then merge children using the tree knapsack approach.

Concretely:
\begin{enumerate}
    \item DFS from root. At each vertex $v$, consider two options:
    \begin{itemize}
        \item Place a sawmill at $v$: cost of $v$'s wood is 0. Children's nearest sawmill is at distance $d_{child}$ (to $v$).
        \item Don't place a sawmill at $v$: cost of $v$'s wood is $w_v \times D$ where $D$ is distance to nearest ancestor sawmill. Children inherit $D + d_{child}$.
    \end{itemize}
    \item Merge children subtrees using knapsack: distribute $j$ sawmills among children.
\end{enumerate}

\subsection{Tree Knapsack Merge}
When vertex $v$ has children $c_1, \ldots, c_m$:
\begin{enumerate}
    \item Start with $g[0] = \text{cost of } v \text{ itself}$.
    \item For each child $c_i$, compute $dp[c_i][j]$ for all $j$, then merge:
    \[
    g'[j_1 + j_2] = \min(g[j_1] + dp[c_i][j_2])
    \]
\end{enumerate}

\section{Complexity Analysis}
\begin{itemize}
    \item \textbf{Time:} $O(N \cdot K^2)$ for the tree knapsack. With the subtree-size bound, this is efficient for $N \leq 100$, $K \leq 50$.
    \item \textbf{Space:} $O(N \cdot K)$.
\end{itemize}

\section{C++ Solution}

\begin{lstlisting}
#include <bits/stdc++.h>
using namespace std;

const long long INF = 1e18;
int N, K;
int w[105], par[105], dist_to_par[105];
vector<int> ch[105]; // children
long long dp[105][55]; // dp[v][j] = min cost, j sawmills in subtree v
int sz[105];

// dist_up[v] = distance from v to its nearest ancestor with a sawmill
// This is passed during DFS

// ancestors: stack of (ancestor_id, distance_from_v_to_ancestor)
// We only need the nearest one's distance.

void dfs(int v, long long distToSawmill) {
    // distToSawmill = distance from v to nearest ancestor with a sawmill
    sz[v] = 1;

    for (int j = 0; j <= K; j++) dp[v][j] = INF;

    // Option 1: don't place sawmill at v
    // Cost for v's wood: w[v] * distToSawmill
    dp[v][0] = (long long)w[v] * distToSawmill;

    // Option 2: place sawmill at v (uses 1 of the j sawmills)
    // Cost for v's wood: 0
    // Children will see d[child] as distance to nearest sawmill
    // We'll handle this by computing dp_with[v][j] separately

    // We need two versions of the subtree DP:
    // dp_no[v][j]: v has no sawmill, ancestor sawmill at distToSawmill
    // dp_yes[v][j]: v has sawmill (counts as 1 of j)

    // Initialize
    vector<long long> dp_no(K + 1, INF), dp_yes(K + 1, INF);
    dp_no[0] = (long long)w[v] * distToSawmill;
    if (K >= 1) dp_yes[1] = 0;

    for (int c : ch[v]) {
        // For dp_no: child c sees distToSawmill + dist_to_par[c]
        // For dp_yes: child c sees dist_to_par[c]

        // Compute child DP under "no sawmill at v" scenario
        dfs(c, distToSawmill + dist_to_par[c]);
        vector<long long> child_no(K + 1);
        for (int j = 0; j <= K; j++) child_no[j] = dp[c][j];

        // Compute child DP under "sawmill at v" scenario
        dfs(c, dist_to_par[c]);
        vector<long long> child_yes(K + 1);
        for (int j = 0; j <= K; j++) child_yes[j] = dp[c][j];

        // Merge dp_no with child_no
        vector<long long> new_no(K + 1, INF);
        for (int j1 = 0; j1 <= K; j1++) {
            if (dp_no[j1] >= INF) continue;
            for (int j2 = 0; j1 + j2 <= K; j2++) {
                if (child_no[j2] >= INF) continue;
                new_no[j1 + j2] = min(new_no[j1 + j2], dp_no[j1] + child_no[j2]);
            }
        }
        dp_no = new_no;

        // Merge dp_yes with child_yes
        vector<long long> new_yes(K + 1, INF);
        for (int j1 = 1; j1 <= K; j1++) {
            if (dp_yes[j1] >= INF) continue;
            for (int j2 = 0; j1 + j2 <= K; j2++) {
                if (child_yes[j2] >= INF) continue;
                new_yes[j1 + j2] = min(new_yes[j1 + j2], dp_yes[j1] + child_yes[j2]);
            }
        }
        dp_yes = new_yes;

        sz[v] += sz[c];
    }

    // Combine: dp[v][j] = min(dp_no[j], dp_yes[j])
    for (int j = 0; j <= K; j++) {
        dp[v][j] = min(dp_no[j], dp_yes[j]);
    }
}

int main(){
    ios::sync_with_stdio(false);
    cin.tie(nullptr);

    cin >> N >> K;

    int root = -1;
    for (int i = 0; i < N; i++) {
        int p;
        cin >> w[i] >> dist_to_par[i] >> p;
        par[i] = p;
        if (p == -1 || p == i) {
            root = i;
        } else {
            ch[p].push_back(i);
        }
    }

    if (root == -1) root = 0;

    // Root has a sawmill (the lumber camp). So root's children see dist_to_par[child].
    // Root's own wood costs 0 (sawmill at root, not counted in K).

    // Process root specially:
    // Root always has sawmill, so we only need dp_yes for root.
    // But to reuse code, call dfs with distToSawmill = 0.
    dfs(root, 0);

    // Answer: dp[root][j] for j <= K, minimum
    long long ans = INF;
    for (int j = 0; j <= K; j++) {
        ans = min(ans, dp[root][j]);
    }
    cout << ans << "\n";

    return 0;
}
\end{lstlisting}

\section{Notes}
The double-DFS approach (calling \texttt{dfs} twice per child) leads to exponential time complexity. The correct approach for the IOI problem is to avoid re-calling DFS by tracking the nearest ancestor sawmill explicitly.

The proper $O(NK^2)$ solution uses:
\[
dp[v][j][a]
\]
where $a$ ranges over ancestors of $v$ (representing which ancestor is the nearest sawmill). Since only ancestors on $v$'s root path matter, and the path has at most $N$ nodes, the total states are bounded. However, for $N \leq 100, K \leq 50$, even simpler approaches work. The key idea remains: tree DP with knapsack merging of children, distinguishing whether the current vertex has a sawmill.

\end{document}
